\chapter{Conclusion and Future Work}
\label{chap:conclusion}

This thesis investigated the integration of Federated Dynamic Averaging (FDA) with Parameter-Efficient Low-Rank Adaptation (LoRA) in federated learning. By formulating and analyzing FDA-RoLoRA---which couples alternating low-rank parameter optimization with an internal dual-track threshold mechanism---this work examined how variance-driven dynamic synchronization can be coupled with alternating low-rank parameter optimization to jointly address spatial parameter overhead and rigid temporal synchronization schedules under statistical data heterogeneity.

\section{Synthesis of Results and Methodological Takeaways}
\label{sec:synthesis_results}

The theoretical formulation, algorithmic design, and empirical evaluation of FDA-RoLoRA yield several core methodological takeaways:

\begin{enumerate}
    \item \textbf{Preservation of Aggregation Linearity}: Dynamic averaging cannot be naively coupled with standard federated LoRA architectures. In simultaneous factor adaptation, bilinear aggregation error---where $\frac{1}{N}\sum A_i B_i \neq (\frac{1}{N}\sum A_i)(\frac{1}{N}\sum B_i)$---distorts local client drift trajectories, causing variance monitoring to misdiagnose true model divergence. By adopting alternating coordinate optimization, the parameter space is restricted to a single factor ($A$ or $B$) per communication round. This guarantees strict linearity in parameter updates, providing the necessary mathematical substrate for stable, decoupled subspace variance tracking without cross-matrix geometric distortion.
    
    \item \textbf{Decoupled Subspace Geometric Monitoring}: Evaluating drift variance directly on isolated low-rank coordinate tensors eliminates the memory and communication overhead of full-model gradient approximations and high-dimensional sketching. By operating an internal dual-track threshold mechanism ($\theta_A, \theta_B$), the system regulates synchronization intervals via a responsive, parameter-free midpoint scaling controller that dynamically calibrates local epoch counts according to realized drift variance on the active coordinate subspace. Establishing the direct Standard update rule as the primary operational controller preserves algorithmic parsimony while maintaining responsive adaptation to local drift.
    
    \item \textbf{Empirical Foundation Model Benchmarks}: Empirical evaluations across transformer foundation models (\texttt{roberta-base} on GLUE tasks) demonstrate that dynamic synchronization achieves substantial communication round savings while preserving model accuracy under statistical data heterogeneity.
\end{enumerate}

\section{Limitations}
\label{sec:limitations}

The methodology and conclusions presented in this work should be considered alongside several operational and experimental limitations:

\begin{enumerate}
    \item \textbf{Empirical Scope and Model Diversity}: While evaluated on \texttt{roberta-base}, extensive multi-seed verification across diverse autoregressive language models and varying non-IID Dirichlet concentration regimes remains valuable future work.
    
    \item \textbf{Simulated Distributed Environment}: Due to the lack of access to the original codebase from the RoLoRA literature, our framework relies on an independent, custom reimplementation. Furthermore, evaluations are conducted within a simulated multi-process federated environment on centralized GPU infrastructure ($2\times$ NVIDIA A10 GPUs), leaving real-world distributed network constraints---such as physical transit latency, asymmetric bandwidth, and stochastic client dropouts---unmodeled.
    
    \item \textbf{Rank Scalability and Alternation Heuristics}: While exact Euclidean variance computation on compact coordinate tensors is computationally negligible for standard low ranks ($r=8$), scaling to substantially higher ranks or wider layer coverage may eventually require approximate sketching techniques (such as the AMS sketch). Additionally, the alternating optimization protocol relies on a rigid deterministic alternation schedule, leaving adaptive or variance-driven matrix toggling criteria unaddressed.
\end{enumerate}

\section{Future Work}
\label{sec:future_work}

Building upon the findings established in this thesis, subsequent research should prioritize several promising directions:

\begin{enumerate}
    \item \textbf{Broader Foundation Model Exploration}: Extending empirical benchmarking to larger autoregressive language models (such as LLaMA and Mistral families) and generative tasks across distributed GPU clusters will further quantify scalability under massive parameter regimes.
    
    \item \textbf{Generalization Across Federated LoRA Variants}: Extending dynamic variance monitoring to alternative parameter-efficient federated paradigms. Key directions include investigating \textbf{FDA + SLoRA} \cite{babakniya2023slora} for coordinate-sparse parameter exchanges, \textbf{FDA + FFA-LoRA} \cite{sun2024improving} with permanently frozen down-projections, and examining whether geometric variance tracking can help diagnose and mitigate client drift in standard simultaneous \textbf{FDA + FedLoRA}. Furthermore, developing adaptive rank allocation algorithms, where the rank $r$ dynamically expands or contracts in response to observed subspace variance, represents a natural extension.
    
    \item \textbf{Advanced Controller Design}: While this thesis adopts the direct, parameter-free Standard update rule, subsequent studies should systematically examine whether advanced closed-loop feedback controllers (such as PID-regulated threshold scaling and variance-acceleration damping) can provide enhanced stability under severe distribution shifts or deeper foundation models. Concurrently, developing formal convergence proofs for alternating dynamic averaging under non-convex loss topologies constitutes an important theoretical milestone.
    
    \item \textbf{Optimization and System Hardening}: From an optimization perspective, exploring asymmetric learning rate schedules tailored independently for down-projection ($A$) and up-projection ($B$) matrices offers potential convergence gains. In addition, architectural integration with low-bit quantization frameworks (such as QLoRA \cite{dettmers2023qlora}) and deployment over physical edge testbeds will further establish operational feasibility. Finally, incorporating FDA-RoLoRA into established open-source federated benchmarking platforms, such as \textbf{FederatedScope-LLM} \cite{kuang2023federatedscope}, will facilitate broader community adoption and standardized evaluation across emerging foundation models.
\end{enumerate}

In conclusion, this work demonstrates that parameter-efficient federated fine-tuning can effectively transcend rigid synchronization schedules through variance-driven dynamic averaging. By establishing that alternating coordinate optimization provides a mathematically consistent substrate for geometric variance monitoring, this thesis provides a principled and communication-efficient framework for federated adaptation.
